\section{Introduction}

Object detection has achieved remarkable progress through the YOLO family of real-time detectors~\cite{yolov8,yolov9,yolo11}.
However, detecting \emph{small objects}---defined as instances occupying fewer than 32$\times$32 pixels---in dense aerial scenes remains significantly more difficult than detecting medium or large targets.
Drone-captured datasets such as VisDrone~\cite{visdrone} exhibit three compounding challenges:
(i)~small objects lack distinctive texture, making them difficult to separate from background noise;
(ii)~standard downsampling in feature pyramids rapidly erodes fine-grained spatial cues needed for precise localization;
and (iii)~class distributions are heavily long-tailed, with \texttt{car} comprising 37\% of instances while \texttt{tricycle} accounts for less than 1\%.

Recent work has introduced targeted modules to address these issues.
Parallel Zone Convolution (Pzconv) expands the receptive field with multi-scale parallel kernels to capture contextual cues surrounding tiny targets.
The Feature Context Module (FCM) employs split-channel cross-attention to suppress background clutter.
RepNCSPELAN4, from the YOLOv9 family, optimizes gradient flow through dense layer aggregation with re-parameterizable branches, preserving shallow high-resolution features critical for small object regression.
While each module addresses part of the problem, no existing design simultaneously handles context expansion, gradient-efficient feature aggregation, attentional purification, and long-tail class rebalancing in a unified framework.

In this work, we propose \textbf{SC-ELAN} (Spatial-Context Efficient Layer Aggregation Network), a modular architecture that unifies these principles.
Our contributions are as follows:

\begin{enumerate}
  \item We introduce \textbf{ContextAwareRepConv}, a multi-branch convolution block that captures multi-scale context (1$\times$1, 3$\times$3, 5$\times$5) during training and re-parameterizes to a single 3$\times$3 convolution at inference, achieving zero additional latency.
  \item We integrate \textbf{LSKA-TSCG} (Large Kernel Spatial Attention with Two-Stage Context Gating) into the SC-ELAN backbone, enabling selective long-range spatial modeling. Through kernel scaling ablation ($k{=}7/11/23$), we confirm that larger kernels consistently improve strict localization metrics.
  \item We design \textbf{DetectCAI}, a training-time class-aware instance reweighting mechanism for the detection head. Through systematic parameter sweeps ($\alpha$, $\beta$), we identify the Soft regime ($\alpha{=}0.10$, $\beta{=}0.40$) as the optimal configuration that improves long-tail categories without degrading dominant classes.
  \item We conduct a \textbf{three-phase systematic ablation} (29 training runs) on VisDrone2019-DET, covering CAI parameter sensitivity, SA-LSKA kernel configurations, and P3-level feature reuse, establishing clear design guidelines for small object detection modules.
\end{enumerate}

Our best model (v2-P1d) achieves mAP50-95 of 0.212 on the VisDrone test-dev set, surpassing all prior SC-ELAN variants while running at $\sim$5.7\,ms per image on an RTX~4090.
